\documentclass[a4paper]{article}
\usepackage{amsfonts}
\usepackage{amsmath}
\usepackage{amssymb}
\usepackage{graphicx}

\begin{document}
	
\noindent \large Auteur: Karanjot Singh\\
\noindent \large Studierichting: Informatica\\
\noindent \large Oefeningenreeks 6A nr. 6.3\\

\normalsize

Gegeven: $x_1 + x_4 = 16$ en $x_3 + x_5 = 28$\\

In een rekenkundige rij geldt: $x_n = x_1 + (n - 1)v$\\

$x_4 = x_1 + 3v \Rightarrow x_1 + (x_1 + 3v) = 16 \Rightarrow 2x_1 + 3v = 16$\\

$x_3 = x_1 + 2v$ en $x_5 = x_1 + 4v \Rightarrow (x_1 + 2v) + (x_1 + 4v) = 28 \Rightarrow 2x_1 + 6v = 28$\\

$2x_1 + 6v - (2x_1 + 3v) = 28 - 16 \Rightarrow 3v = 12 \Rightarrow v = 4$\\

$2x_1 + 12 = 16 \Rightarrow 2x_1 = 4 \Rightarrow x_1 = 2$\\

Dus:\\

$x_1 = 2$\\

$x_2 = x_1 + v = 6$\\

$x_3 = x_1 + 2v = 10$\\

\begin{thebibliography}{999}
	%\bibitem{a} Referentie
\end{thebibliography}

\end{document}
